\chapter{Projeto 1: O Jogo da Coleta}

\begin{objetivos}
\begin{itemize}
\item Juntar o que você aprendeu em um jogo completo;
\item Usar clones, variáveis, condições e eventos juntos;
\item Criar pontos, vidas e uma tela de fim de jogo.
\end{itemize}
\end{objetivos}

\section{O plano do jogo}

Neste projeto você vai criar um jogo em que \textbf{maçãs caem do céu} e você controla uma \textbf{cesta} para pegá-las. Antes de começar, o roteiro (lembra do planejamento?):

\begin{definicao}[Roteiro do jogo]
\textbf{Atores:} Cesta e Maçã.\\
\textbf{Variáveis:} \textbf{pontos} e \textbf{vidas}.\\
\textbf{Regras:} pegar a maçã dá 1 ponto; deixar a maçã cair no chão tira 1 vida; sem vidas, o jogo acaba.
\end{definicao}

\begin{missao}[Missão 17: Monte o Jogo da Coleta]
Siga os cinco passos e, no final, você terá o seu primeiro jogo completo!
\end{missao}

\section{Passo 1: a cesta}

Escolha ou desenhe uma cesta. Coloque-a embaixo do palco e faça-a andar com as setas.

\begin{codigo}[Cesta]
\begin{pilha}
\bk[scEvents]{quando \bandeira\ for clicado}
\bk[scMotion]{vá para x: \entrada{0} y: \entrada{$-150$}}
\end{pilha}
\begin{pilha}
\bk[scEvents]{quando a tecla \menu{seta direita} for pressionada}
\bk[scMotion]{mude x em \entrada{15}}
\end{pilha}
\begin{pilha}
\bk[scEvents]{quando a tecla \menu{seta esquerda} for pressionada}
\bk[scMotion]{mude x em \entrada{$-15$}}
\end{pilha}
\end{codigo}

\section{Passo 2: as variáveis}

Crie \textbf{pontos} e \textbf{vidas} (para todos os atores). Dê ao Palco dois cenários, \textbf{Jogo} e \textbf{Fim de jogo}. Ao começar, volte ao cenário Jogo, zere os pontos e dê 3 vidas.

\begin{codigo}[No Palco]
\begin{pilha}
\bk[scEvents]{quando \bandeira\ for clicado}
\bk[scLooks]{mude o cenário para \menu{Jogo}}
\bk[scVariables]{defina \menu{pontos} para \entrada{0}}
\bk[scVariables]{defina \menu{vidas} para \entrada{3}}
\end{pilha}
\end{codigo}

\section{Passo 3: a maçã e os clones}

O ator Maçã fica escondido e vai criando clones em lugares aleatórios do topo.

\begin{codigo}[Maçã original]
\begin{pilha}
\bk[scEvents]{quando \bandeira\ for clicado}
\bk[scLooks]{esconda}
\sbc[scControl]{sempre}
\bk[scMotion]{vá para x: \reporter{número aleatório entre \entrada{$-200$} e \entrada{200}} y: \entrada{170}}
\bk[scControl]{crie um clone de \menu{mim mesmo}}
\bk[scControl]{espere \entrada{1} segundos}
\sbcend
\end{pilha}
\end{codigo}

\section{Passo 4: o que a maçã faz}

Cada clone cai. Se tocar na cesta, dá 1 ponto. Se tocar no chão, tira 1 vida.

\begin{codigo}[Clone da maçã]
\begin{pilha}
\bk[scControl]{quando eu começar como um clone}
\bk[scLooks]{mostre}
\sbc[scControl]{repita até que \booleano{\variavel{y} $<$ \entrada{$-160$}}}
\bk[scMotion]{mude y em \entrada{$-5$}}
\sbc[scControl]{se \booleano{tocando em \menu{Cesta}?} então}
\bk[scVariables]{mude \menu{pontos} em \entrada{1}}
\bk[scControl]{apague este clone}
\sbcend
\sbcend
\bk[scVariables]{mude \menu{vidas} em \entrada{$-1$}}
\bk[scControl]{apague este clone}
\end{pilha}
\end{codigo}

\begin{dica}
\variavel{y} é um bloco da categoria \bloco[scMotion]{Movimento} que mostra a posição vertical do ator. Quando ela fica menor que $-160$, a maçã chegou ao chão.
\end{dica}

\section{Passo 5: fim de jogo}

Quando as vidas acabam, o Palco avisa que o jogo terminou e para tudo.

\begin{codigo}[No Palco]
\begin{pilha}
\bk[scEvents]{quando \bandeira\ for clicado}
\sbc[scControl]{repita até que \booleano{\variavel{vidas} $<$ \entrada{1}}}
\bk[scControl]{espere \entrada{0.1} segundos}
\sbcend
\bk[scLooks]{mude o cenário para \menu{Fim de jogo}}
\bk[scControl]{pare \menu{tudo}}
\end{pilha}
\end{codigo}

\begin{cuidado}
Se as maçãs não caírem, confira se o ator Maçã tem o bloco \bloco[scLooks]{esconda} só no \textbf{original}, e o bloco \bloco[scLooks]{mostre} no começo do \textbf{clone}. Se as pontuações não subirem, veja se o nome do ator no bloco \bloco[scSensing]{tocando em} está certo.
\end{cuidado}

\section{Melhorando o jogo}

\begin{exercicio}[Maçã mais rápida a cada ponto]
Faça a maçã cair mais rápido conforme os pontos aumentam. Dica: use \variavel{pontos} no valor de queda.
\end{exercicio}

\begin{solucao}
Troque o \bloco[scMotion]{mude y em \entrada{$-5$}} do clone por um valor que dependa dos pontos:
\begin{pilha}
\bk[scMotion]{mude y em \reporter{\entrada{$-3$} $-$ \reporter{\variavel{pontos} / \entrada{5}}}}
\end{pilha}
A cada 5 pontos, a maçã cai mais 1 passo por volta.
\end{solucao}

\begin{exercicio}[Maçã podre]
Adicione uma segunda fantasia ``podre''. Quando a cesta pegar a podre, perde 1 vida em vez de ganhar ponto.
\end{exercicio}

\begin{solucao}
Ao criar o clone, sorteie a fantasia:
\begin{pilha}
\bk[scControl]{quando eu começar como um clone}
\bk[scLooks]{mude para a fantasia \reporter{número aleatório entre \entrada{1} e \entrada{2}}}
\bk[scLooks]{mostre}
\end{pilha}
E na hora de tocar na cesta, veja qual fantasia o clone tem: \bloco[scControl]{se \booleano{\variavel{número da fantasia} = \entrada{2}} então} tira vida; \bloco[scControl]{senão} dá ponto.
\end{solucao}

\begin{exercicio}[Som de ponto]
Toque um som toda vez que a cesta pegar uma maçã.
\end{exercicio}

\begin{solucao}
Dentro do \bloco[scControl]{se \booleano{tocando em \menu{Cesta}?} então}, antes de apagar o clone:
\begin{pilha}
\bk[scSound]{inicie o som \menu{Pop}}
\end{pilha}
\end{solucao}

\begin{desafio}[Sua versão]
Troque a maçã por outra coisa (estrelas, moedas, peixes\ldots) e a cesta por um personagem seu. Acrescente um cronômetro de 30 segundos. A partida termina quando o tempo ou as vidas acabam, o que vier primeiro. A pontuação fica visível no final.
\end{desafio}

\begin{solucao}
Mantenha a pilha que verifica as vidas. Marque a caixinha da variável \textbf{pontos} na categoria Variáveis para deixar o placar visível. Crie \textbf{tempo} e acrescente esta pilha ao Palco:
\begin{pilha}
\bk[scEvents]{quando \bandeira\ for clicado}
\bk[scVariables]{defina \menu{tempo} para \entrada{30}}
\sbc[scControl]{repita até que \booleano{\variavel{tempo} = \entrada{0}}}
\bk[scControl]{espere \entrada{1} segundos}
\bk[scVariables]{mude \menu{tempo} em \entrada{$-1$}}
\sbcend
\bk[scLooks]{mude o cenário para \menu{Fim de jogo}}
\bk[scControl]{pare \menu{tudo}}
\end{pilha}
\end{solucao}
